\section{Development Data and Protocol}
\label{app:development-protocol}
\subsection{Data and sampling}
The development evaluation uses ContractNLI's annotated contract--hypothesis decisions \citep{koreeda2021contractnli}. We sort development documents by integer ID and sample 30 without replacement using seed 20260920. Each example retains the full contract text, all 17 original hypotheses, and the original class annotations. Evidence spans are excluded from model inputs, and sampling is independent of model predictions.

The sampled 510 targets comprise 280 entailments, 48 contradictions, and 182 not-mentioned labels, motivating class-level reporting alongside overall accuracy. Each experimental condition partitions the same hypotheses into request groups. We match predictions by document and hypothesis before paired analysis. The five-contract panel uses the first five documents in the randomized sample.

We independently recompute scores and transitions from individual predictions, verifying that paired observations share the same targets and gold labels. All three pilot phases have complete prediction coverage; Appendix~\ref{app:repro} gives the document selection and sample sizes. Scores measure agreement with the original annotations, whose legal interpretations may remain ambiguous in individual cases.

\subsection{Models and inference settings}
The development baseline compares Jev 1.13.0, Qwen3.5-4B and 9B, Gemini 3.1 Pro Preview, Gemini 3.5 Flash-Lite, GPT-5.6 Luna and Terra, GPT-6 Astra, Claude Sonnet 5, and Claude Haiku 4.5. Jev receives the contract as shared state and each hypothesis as a native Choice question. Generative models receive structured classification instructions and return a JSON label map. Task instructions are semantically aligned across these model-specific interfaces.

In the initial baseline, grouping, and anchor experiments, the Qwen models use BF16, disabled thinking, and temperature zero. Gemini Pro uses low thinking, Flash-Lite minimal thinking, Luna no reasoning, and Astra high reasoning. Luna and Astra are accessed through OpenRouter, with OpenAI identified as the provider for the initial pilot. Reasoning budgets differ across models, and actual reasoning length can vary within a setting: all 30 Astra baseline responses contain reasoning tokens, totaling 23,042.

The later development extension adds Terra, Sonnet, and Haiku with the same inference settings used in their test evaluations: medium reasoning for Terra, adaptive thinking with medium effort for Sonnet, and disabled thinking for Haiku. Terra and Sonnet permit 16,384 output tokens; Haiku permits 8,192. All three use the Convert gateway. Gemini Pro's extended controls retain its original low-thinking configuration. The new GPT-6 controls use high reasoning through Convert, which identifies the model as GPT-6; the earlier baseline uses GPT-6 Astra through OpenRouter. We keep these controls separate because the two routes are not assumed to serve an identical checkpoint.

\subsection{Experimental conditions}
The initial baseline evaluates seven models on the same 30 contracts, with one $K=17$ request per contract. Ninety observations for Jev, Gemini Pro, and Qwen4B come from an earlier collection, followed by evaluation of the other four models. Collection times therefore differ. For the earlier observations, request duration serves as a proxy for document completion time where the latter was not measured separately.

All seven models are then evaluated under seven conditions on the common five-contract subset. The same conditions are applied to Jev, the two Qwen models, Flash-Lite, and Luna on all 30 contracts. This expansion reuses the development sample after exploratory inspection. Each sensitivity comparison uses the $K=17$ baseline from its own experimental phase.

For maximum batch sizes $K\in\{1,4,8,17\}$, the 17 judgments produce 17, 5, 3, and 1 requests respectively, including the final partial batch. The grouping comparison preserves hypothesis order. A permutation control uses a document-specific shuffle determined by seed 20260920 and the document ID. One additional $K=17$ request repeats the original input. Grouping, permutation, and repeat conditions use concurrency one; an additional $K=1$ condition permits four simultaneous requests. The order of model--condition groups is randomized separately in each evaluation. Evaluations may overlap in time, and prefix caching is enabled for local models.

\paragraph{Extended development comparisons.}
After inspecting the earlier development and test results, we extend the development comparison using the same contracts, prompts, and labels. Sonnet, Terra, and Haiku each receive the $K=17$ baseline, one unchanged repeat, and one permutation on all 30 contracts. Each model's baseline supports both the cost--accuracy comparison and its paired controls. Gemini Pro and GPT-6 receive the same three conditions in separate new collections. Their controls use these newly matched baselines, while Figure~\ref{fig:development-map} retains the earlier Gemini Pro and OpenRouter GPT-6 Astra baselines. Every comparison pairs responses within one collection and route. The extension contains 450 requests for these five panels and 1,080 requests for the three new anchor evaluations, with no smaller-group or concurrency conditions. Model selection remains exploratory; samples, inference settings, and scoring rules are fixed before these additional requests.

\paragraph{Uncertainty, latency, and cost.}
Dependence among NLP observations matters for statistical evaluation~\citep{dror2018hitchhiker}. We report descriptive paired percentile intervals from 5,000 bootstrap resamples of whole contracts, preserving within-contract dependence, without multiple-comparison correction. Zero observed changes can yield a degenerate interval despite uncertainty about the underlying change probability. Latency is the median elapsed time required to obtain all requested decisions for a document, including concurrent execution where applicable. API costs use provider-reported charges or recorded token usage priced at public rates. These timings combine model inference, network latency, scheduling, caching, and provider infrastructure.

For local models, we estimate rental-equivalent inference cost per contract as
\begin{equation}
\widehat C_{\mathrm{GPU}} = \frac{r}{30\times3600}\sum_{d=1}^{30} t_d,
\end{equation}
where $t_d$ is the observed elapsed time in seconds for contract $d$ in the same 30-contract, $K=17$, concurrency-one baseline, and $r$ is \$2.00 per GPU-hour. Each model uses one GPU. This rounded reference lies within three published RTX PRO 6000 rental rates of \$1.80--\$2.09 per hour; it is not a market average.\footnote{Rates accessed September 21, 2026: \href{https://nebius.com/prices}{Nebius}, \$1.80; \href{https://verda.com/pricing}{Verda}, \$1.86; and \href{https://www.runpod.io/pricing}{Runpod}, \$2.09 per GPU-hour.} The estimates value measured local time rather than actual GPU charges. The local Max-Q Workstation GPU differs from the cloud configurations, which include Server Edition hardware; equal cloud throughput therefore cannot be inferred. Model loading, preliminary checks, idle allocation, storage, transfer, and taxes are excluded.

\begin{table*}[t]
\centering
\small
\begin{tabularx}{\textwidth}{Xrrrrr}
\toprule
\rowcolor{tableHeader}
Model & \shortstack{Accuracy\\(\%) $\uparrow$} & \shortstack{Macro-F1\\(\%) $\uparrow$} & Seconds $\downarrow$ & \shortstack{USD/contract\\$\downarrow$} & Valid calls \\
\midrule
\rowcolor{jevRow}
\modelicon{jev}Jev 1.13.0 & 77.06 & 70.72 & \bestscore{1.25} & \bestscore{0.000250} & 30/30 \\
\addlinespace[3pt]
\modelicon{gemini}Gemini 3.5 Flash-Lite & 80.98 & 74.02 & 1.79 & 0.001486 & 30/30 \\
\modelicon{gemini}Gemini 3.1 Pro Preview & \bestscore{82.94} & \bestscore{77.19} & 3.79 & 0.009502 & 30/30 \\
\addlinespace[3pt]
\modelicon{openai}GPT-5.6 Luna & 77.65 & 70.23 & 1.64 & 0.001058 & 30/30 \\
\modelicon{openai}GPT-5.6 Terra & \secondscore{81.76} & \secondscore{75.90} & 11.47 & 0.011318 & 30/30 \\
\modelicon{openai}GPT-6 Astra & \secondscore{81.76} & 74.64 & 20.35 & 0.089897 & 30/30 \\
\addlinespace[3pt]
\modelicon{claude}Claude Haiku 4.5 & 78.82 & 73.46 & 5.73 & 0.005906 & 30/30 \\
\modelicon{claude}Claude Sonnet 5 & 81.37 & 75.09 & 2.75 & 0.012868 & 30/30 \\
\addlinespace[3pt]
\modelicon{qwen}Qwen3.5-4B & 78.43 & 66.53 & \secondscore{1.38} & \secondscore{0.000783} & 30/30 \\
\modelicon{qwen}Qwen3.5-9B & 78.63 & 73.15 & 2.89 & 0.001627 & 30/30 \\
\bottomrule
\end{tabularx}
\caption{Development baseline on 30 contracts and 510 labels, requesting all 17 judgments together. The original seven baseline results are retained; Sonnet, Terra, and Haiku are evaluated on the same contracts. Seconds denotes median document latency. API costs use reported charges or published token rates; GPU costs follow the development calculation in Appendix~\ref{app:development-protocol}. Hardware, API routing, and inference settings differ across models. $\uparrow$/$\downarrow$: higher/lower is better. \textbf{Bold} and \underline{underlining} indicate the best and second-best distinct values, including ties, within each indicated comparison; these marks do not imply statistical significance.}
\label{tab:baseline}
\end{table*}

